% ============================================================
%  Chapter 1: The Ion's Journey --- Classical and Quantum in Concert
%
%  Source: union/publication/crowns/ion-journey.tex
%  This chapter adapts that paper into book-chapter format.
%  The ion: m/z = 299.0555, traced through four platforms.
% ============================================================

\epigraph{%
  There is no experiment that distinguishes the two descriptions.\\
  There is only the experiment.%
}{}

%% ── §1.1 ─────────────────────────────────────────────────────────────────────
\section{Prologue: One Molecule, Many Languages}
\label{sec:ch1-prologue}

Somewhere in a solution of buffer and analyte, a molecule acquires a proton
and becomes an ion.
The molecule --- let us call it M --- has monoisotopic mass $299.0555\,\mathrm{Da}$
and carries a single positive charge.
Its $m/z$ ratio is therefore $299.0555$.
In the next few microseconds it will be accelerated by electric fields,
focused by RF multipoles, separated according to its mass-to-charge ratio,
fragmented, separated again, and finally detected.
At the end of this journey a number will be written down: the measured $m/z$.

We will describe this journey twice.

The first description uses classical mechanics: Newton's equations of motion,
the Lorentz force, viscous drag, and conservation of energy.
Every ion is a charged particle with a definite position and velocity at every
instant.
The second description uses quantum mechanics: a wave function that evolves
under a Hamiltonian, collapses on interaction with the instrument, and produces
eigenvalues that are the observable quantities.

Both descriptions are standard.
Neither is approximate.
Both give the same measured $m/z$ to within two parts per million across four
different mass analyser technologies.

The agreement is not a coincidence.
Both descriptions are consequences of the same fact about the ion: its phase
space is bounded.

%% ── §1.2 ─────────────────────────────────────────────────────────────────────
\section{Stage One: Electrospray Ionisation}
\label{sec:ch1-esi}

The ion is born in an electrospray droplet.

\subsection*{Classical description}
A liquid jet of analyte solution is driven through a capillary held at several
kilovolts above the inlet of the mass spectrometer.
The electric field at the capillary tip overcomes the surface tension of the
liquid; the jet breaks into a spray of highly charged droplets.
Each droplet carries a net positive charge and begins to evaporate.
The Rayleigh instability criterion --- charge repulsion exceeds surface tension
when
\begin{equation}\label{eq:rayleigh}
  q > 8\pi \sqrt{\varepsilon_0 \gamma R^3},
\end{equation}
where $\gamma$ is surface tension and $R$ is droplet radius --- determines
when the droplet undergoes a Coulomb explosion, emitting a stream of smaller
droplets and, ultimately, individual ions.
The ion M$+$H$^+$ emerges with a velocity distribution set by the thermal
energy of the solution, $k_\mathrm{B} T \approx 25\,\mathrm{meV}$ at room
temperature.

\subsection*{Quantum mechanical description}
The protonation event is the collapse of the molecule's quantum state onto
a particular eigenstate of the Hamiltonian $\hat{H}(\mathbf{r}_\mathrm{H^+})$
describing the proton--molecule interaction.
The resulting wave function $\psi(\mathbf{r}, t)$ of the ion factorises,
in the Born--Oppenheimer approximation, into an electronic part
$\phi_\mathrm{el}(\mathbf{r}_\mathrm{el})$ and a nuclear part
$\chi(\mathbf{R}_\mathrm{nuc}, t)$.
The nuclear wave function is a Gaussian wave packet whose centre follows
the classical Newtonian trajectory and whose
width is $\Delta x \approx \hbar / (m_\mathrm{ion} v_0)^{1/2}$, negligibly
small for a 299-Da ion.

\subsection*{Why they agree}
The ion's phase space volume immediately after ionisation is
\begin{equation}
  V_\Gamma = \Delta x^3\, \Delta p^3 \sim \left(\frac{\hbar}{m v_0}\right)^{3/2}
             \cdot (m \cdot k_\mathrm{B}T)^{3/2} = (\hbar)^{3/2}(k_\mathrm{B}T)^{3/2},
\end{equation}
which is bounded.
The ion's thermal de Broglie wavelength $\lambda = h/(2\pi m k_\mathrm{B}T)^{1/2}$
is of order $10^{-13}\,\mathrm{m}$ for a 299-Da molecule at 300\,K ---
sixteen orders of magnitude below the spatial resolution of any mass
spectrometer.
The two descriptions are therefore identical in every measurable respect:
the wave function is a point in phase space that satisfies Newton's equations.

\begin{remark}
  The equivalence here is not the correspondence principle.
  It is not that the quantum number $n$ is large.
  It is that the ion's phase space volume is a finite, bounded quantity, and
  therefore both the classical trajectory (which samples phase space
  continuously) and the quantum state vector (which is a density on phase
  space) describe the same bounded region.
  The Bounded Phase Space Law, proved in Part~II, guarantees this equivalence
  for any persistent physical system.
\end{remark}

%% ── §1.3 ─────────────────────────────────────────────────────────────────────
\section{Stage Two: Chromatographic Separation and Ion Transfer}
\label{sec:ch1-chrom}

After electrospray, the ions pass through a series of RF multipoles and
ion funnels before reaching the mass analyser.

\subsection*{Classical description}
In a quadrupole ion guide, the ion experiences an RF electric field
\begin{equation}
  \mathbf{E}(x,y,t) = \frac{V_0 \cos(\Omega t)}{r_0^2}(x\,\hat{x} - y\,\hat{y}),
\end{equation}
where $V_0$ is the RF voltage amplitude, $\Omega$ the RF angular frequency,
and $r_0$ the inscribed radius of the quadrupole.
The equations of motion reduce to the Mathieu equation
\begin{equation}
  \ddot{u} + [a_u - 2q_u \cos(2\xi)]\,u = 0,
  \qquad \xi = \tfrac{1}{2}\Omega t,
\end{equation}
with stability parameters $a_u$ and $q_u$ depending on DC and RF voltages.
Stable trajectories are bounded; unstable trajectories diverge exponentially
and the ion is lost.

\subsection*{Quantum mechanical description}
In the quantum picture the RF multipole applies a periodic potential to the
ion's wave function.
Floquet theory decomposes the time-periodic Hamiltonian into quasi-energy bands.
The condition for a stable, persistent ion trajectory is that the quasi-energy
lies within a band gap --- equivalently, the ion's Floquet state is bounded.
The stable Floquet states are precisely those whose $a_u, q_u$ parameters
fall within the stability diagram of the Mathieu equation.

\subsection*{Why they agree}
The Mathieu stability diagram is simultaneously the classical stability
criterion and the quantum quasi-energy gap condition.
Both descriptions produce the same boundary in $(a_u, q_u)$ space, because
both are asking the same question: is the ion's phase space trajectory bounded?

%% ── §1.4 ─────────────────────────────────────────────────────────────────────
\section{Stage Three: Mass Analysis on Four Platforms}
\label{sec:ch1-analyzers}

The ion is now measured four times: in a time-of-flight (TOF) analyser,
an Orbitrap, a Fourier-transform ion cyclotron resonance (FT-ICR) cell,
and a quadrupole mass filter.
Each platform operates on a different physical principle.
All four return $m/z = 299.0555 \pm 0.0006$ (2\,ppm).

\subsection{Time-of-Flight Analysis}
\label{sec:ch1-tof}

\paragraph{Classical.}
The ion is accelerated through a potential $U$ and enters a field-free flight
tube of length $L$.
Its velocity is $v = \sqrt{2eU/m}$, and its flight time is
\begin{equation}\label{eq:tof-classical}
  t = \frac{L}{v} = L\sqrt{\frac{m}{2eU}}.
\end{equation}
Measuring $t$ to nanosecond precision with a known $L$ and $U$ gives $m$ to
sub-ppm accuracy.

\paragraph{Quantum.}
The ion wave packet entering the flight tube has central momentum
$p_0 = m v_0 = \sqrt{2meU}$.
The wave packet propagates freely: the phase accumulated in traversing
length $L$ is
\begin{equation}\label{eq:tof-quantum}
  \phi = \frac{p_0 L}{\hbar} = \frac{L\sqrt{2meU}}{\hbar}.
\end{equation}
The time of arrival is determined by the group velocity $\partial\omega/\partial k
= p_0/m$, giving the same $t$ as Eq.~\eqref{eq:tof-classical}.

\paragraph{Agreement.}
Equations~\eqref{eq:tof-classical} and~\eqref{eq:tof-quantum} are identical
to sub-femtosecond precision.
The quantum correction (spread of the wave packet) is of order
$\hbar/(p_0 L) \sim 10^{-20}$, unmeasurable.

\subsection{Orbitrap Analysis}
\label{sec:ch1-orbitrap}

\paragraph{Classical.}
Inside an Orbitrap, the ion orbits a central spindle electrode in a
hyper-logarithmic electrostatic field.
The axial equation of motion is that of a simple harmonic oscillator:
\begin{equation}\label{eq:orbitrap-axial}
  \ddot{z} + \omega_z^2 z = 0, \qquad
  \omega_z = \sqrt{\frac{eC}{m}},
\end{equation}
where $C$ is an instrument constant.
The axial oscillation frequency $\omega_z$ is measured by Fourier-transforming
the image current induced on the outer electrodes, giving $m/e = C/\omega_z^2$.

\paragraph{Quantum.}
In quantum mechanics, the Orbitrap potential is a harmonic well.
The energy eigenvalues are $E_n = \hbar\omega_z(n + \tfrac{1}{2})$.
The ion in a coherent state (the quantum state that most closely mimics a
classical oscillator) oscillates at exactly the classical frequency $\omega_z$,
with expectation value $\langle z \rangle = A\cos(\omega_z t)$.
The induced image current is therefore identical to the classical prediction.

\paragraph{Agreement.}
The Orbitrap is, literally, a quantum harmonic oscillator in a coherent state.
The classical frequency and the quantum transition frequency are the same:
\begin{equation}
  \omega_z = \sqrt{\frac{eC}{m}} \quad \text{(classical)} \qquad = \qquad
  \frac{E_{n+1} - E_n}{\hbar} \quad \text{(quantum)}.
\end{equation}
This is not approximate; it is exact.

\subsection{FT-ICR Analysis}
\label{sec:ch1-fticr}

\paragraph{Classical.}
In a magnetic field $\mathbf{B} = B\hat{z}$, the ion undergoes circular
cyclotron motion at frequency
\begin{equation}\label{eq:cyclotron}
  \omega_c = \frac{eB}{m}.
\end{equation}
A chirp RF excitation pulse brings all ions of a given mass to a large cyclotron
radius; the dephasing of the resulting image current, Fourier-transformed, yields
$\omega_c$ and hence $m$.

\paragraph{Quantum.}
The Landau levels in a magnetic field are
$E_n = \hbar\omega_c(n + \tfrac{1}{2})$.
The spacing is $\Delta E = \hbar\omega_c$.
The RF pulse drives a coherent superposition of Landau levels; the expectation
value $\langle x + \iu y\rangle$ rotates at $\omega_c$.
The image current observed at the electrodes is identical to the classical
cyclotron signal.

\paragraph{Agreement.}
The cyclotron frequency is the Larmor frequency of the Landau-level ladder.
Classical orbit and quantum eigenvalue spacing agree exactly:
$\omega_c = eB/m$ in both languages.

\subsection{Quadrupole Mass Filter}
\label{sec:ch1-quad}

\paragraph{Classical.}
A quadrupole mass filter operates at the apex of the Mathieu stability diagram:
for a narrow range of $a_u, q_u$ around the stability apex, only ions of a
specific $m/z$ have bounded trajectories and are transmitted.
Scanning the DC/RF ratio scans the transmitted mass.

\paragraph{Quantum.}
The quasi-energy bands of the periodic Hamiltonian contain a narrow allowed
band at the stability apex.
The bandpass condition is $|\Delta E_\mathrm{Floquet}| < \delta$, where
$\delta$ is the stability width.
The transmitted ions are those in the allowed quasi-energy band, which
corresponds identically to the classical stability apex.

\paragraph{Agreement.}
The quadrupole filter is a Floquet bandpass filter.
Classical stability and quantum allowed band are the same mathematical object:
the bounded region of Mathieu stability space.

\subsection{The Four-Platform Summary}
\label{sec:ch1-fourplatform}

\begin{table}[htbp]
\centering
\caption{%
  The same ion ($\mz = 299.0555$) measured on four platforms.
  Classical and quantum descriptions give identical results; both agree with
  experiment to $\leq 2\,\mathrm{ppm}$.
}
\label{tab:four-platforms}
\begin{tabular}{@{}llll@{}}
\toprule
Platform & Physical principle & Classical formula & Quantum analogue \\
\midrule
TOF      & Flight time        & $t = L\sqrt{m/2eU}$ & Group velocity of wave packet \\
Orbitrap & Axial oscillation  & $\omega_z = \sqrt{eC/m}$ & Coherent state of harmonic oscillator \\
FT-ICR   & Cyclotron orbit    & $\omega_c = eB/m$ & Landau level spacing \\
Quadrupole & Mathieu stability & Apex of stability diagram & Floquet quasi-energy band \\
\bottomrule
\end{tabular}
\end{table}

In every case the classical formula and the quantum formula are not
approximately equal --- they are exactly equal.
The ion's phase space is bounded; therefore, classical trajectory and quantum
wave function describe the same bounded motion in the same bounded region.

%% ── §1.5 ─────────────────────────────────────────────────────────────────────
\section{Stage Four: Fragmentation}
\label{sec:ch1-frag}

Having been measured, the ion is isolated and subjected to collision-induced
dissociation (CID).
It collides with nitrogen gas molecules and absorbs internal energy sufficient
to break one or more bonds.

\subsection*{Classical description}
Each collision transfers kinetic energy to internal vibrational modes.
The energy redistribution follows RRKM theory: the rate constant for bond
cleavage is
\begin{equation}
  k(E) = \frac{N^\ddagger(E - E_0)}{h\,\rho(E)},
\end{equation}
where $N^\ddagger$ is the number of states at the transition state above
threshold $E_0$, and $\rho(E)$ is the density of states at total internal
energy $E$.
Fragmentation is a classical transition-state theory calculation.

\subsection*{Quantum mechanical description}
The internal energy of the ion is quantised: $E = \sum_i n_i \hbar\omega_i$
where $\omega_i$ are the normal-mode frequencies and $n_i$ are occupation
numbers.
The fragmentation rate is the quantum mechanical tunnelling rate through the
potential energy surface barrier, corrected by the density of states --- which
is the quantum origin of the RRKM formula.
The classical RRKM result is the leading term of the quantum expression
in the limit of high density of states.

\subsection*{Why they agree}
For a 299-Da molecule with dozens of vibrational modes, the density of states
at CID energies (1--10\,eV) is $\rho(E) \sim 10^{12}\,\mathrm{states/eV}$.
The quantum level spacing is $\sim 10^{-12}\,\mathrm{eV}$, far below any
experimental energy resolution.
Classical RRKM and quantum rate theory are numerically identical because the
ion's vibrational phase space, though discrete, is so densely packed that the
classical continuum approximation is exact.
The bounded partition space of the ion's vibrational degrees of freedom is
large enough that the classical limit applies; the Bounded Phase Space Law
guarantees its convergence.

%% ── §1.6 ─────────────────────────────────────────────────────────────────────
\section{Stage Five: Detection}
\label{sec:ch1-detection}

Fragment ions strike a detector --- either a microchannel plate (MCP) or the
induction detector of the Orbitrap/FT-ICR.

\subsection*{Classical description}
An ion of charge $e$ and velocity $v$ impinging on an MCP generates a
cascade of secondary electrons.
The gain is classical: $G \approx \exp(\alpha \ell)$ where $\alpha$ is the
secondary-emission coefficient and $\ell$ the channel length.
The total induced charge is $Q = Ge$.

\subsection*{Quantum mechanical description}
Each detected ion is a single quantum event: a discrete charge $e$ arriving at
the detector.
The cascade process is governed by quantum mechanical secondary emission
probabilities, but in the limit of large gain $G \gg 1$, the output distribution
is Poissonian and indistinguishable from the classical result.

\subsection*{Agreement}
The detector registers a number of electrons.
Whether one views each electron as a classical charge or as the outcome of a
quantum measurement makes no difference to the final mass spectrum.

%% ── §1.7 ─────────────────────────────────────────────────────────────────────
\section{Epilogue: What the Ion Has Shown}
\label{sec:ch1-epilogue}

We have now traced a single ion from its birth in the electrospray droplet to
its death on the detector.
At each of the five stages, we wrote down two descriptions --- classical and
quantum --- and found that they agreed to better precision than any instrument
can distinguish.

The agreement is not circumstantial.
It does not depend on the ion being large (a 299-Da molecule is small by
biochemical standards).
It does not depend on the approximation that quantum mechanics ``reduces'' to
classical mechanics at large quantum numbers.
It holds because the ion's phase space is bounded.

A bounded phase space has three properties that force the two descriptions into
agreement:
\begin{enumerate}
  \item \textbf{Finite mode spectrum.}
    A bounded oscillator has a countable set of frequencies.
    Quantum eigenvalues and classical Fourier modes are the same discrete set.
  \item \textbf{Poincaré recurrence.}
    Every trajectory returns arbitrarily close to its initial condition.
    The classical trajectory is periodic; the quantum wave function is
    quasi-periodic.
    Both sample the same bounded region of phase space.
  \item \textbf{Hierarchical partitionability.}
    The bounded phase space can be partitioned into cells of decreasing size.
    A classical measurement picks a cell; a quantum measurement collapses to an
    eigenstate.
    For a 299-Da ion the cells are so small ($\Delta x \sim 10^{-13}$\,m)
    that no instrument can distinguish between picking a cell and collapsing
    to an eigenstate.
\end{enumerate}

These three properties are the three clauses of the Bounded Phase Space Law.
They are the subject of Part~II.

\bigskip
The ion has demonstrated something profound.
Two crowns of physics --- one classical, one quantum --- have descended on the
same head.
The rest of this monograph is an account of why they fit.

\vspace{2ex}
\begin{tcolorbox}[colback=crownblue!5, colframe=crownblue!60!black,
                  title={\textbf{Chapter 1 Summary}}]
\begin{itemize}[noitemsep]
  \item A single ion ($m/z = 299.0555$) was traced through five stages of a
    mass spectrometric measurement.
  \item At each stage, the classical and quantum descriptions were written down
    and shown to agree to within instrumental precision ($\leq 2\,\mathrm{ppm}$).
  \item The agreement holds because the ion's phase space is bounded.
    A bounded phase space forces: finite mode spectrum; Poincaré recurrence;
    hierarchical partitionability.
  \item These three properties are the clauses of the Bounded Phase Space Law,
    derived in Chapter~\ref{ch:axiom}.
\end{itemize}
\end{tcolorbox}
